\documentclass[11pt]{article}

\usepackage[T1]{fontenc}
\usepackage{amsmath,amssymb,amsthm}
\usepackage{hyperref}

\title{Higher-Order Asymptotics of the Discrepancy Function\\
under Fixed Model and Distributional Misspecification}
\author{}
\date{\today}

\newcommand{\E}{\mathbb{E}}
\newcommand{\tr}{\operatorname{tr}}
\newcommand{\Var}{\operatorname{Var}}
\newcommand{\Cov}{\operatorname{Cov}}
\newcommand{\vech}{\operatorname{vech}}
\newcommand{\rank}{\operatorname{rank}}
\newcommand{\argmin}{\operatorname*{arg\,min}}
\newcommand{\argmax}{\operatorname*{arg\,max}}
\newcommand{\plim}{\operatorname*{plim}}
\newcommand{\rmsea}{\widehat{\varepsilon}}
\newcommand{\dto}{\xrightarrow{d}}
\newcommand{\pto}{\xrightarrow{p}}

\theoremstyle{plain}
\newtheorem{theorem}{Theorem}
\newtheorem{lemma}{Lemma}
\newtheorem{corollary}{Corollary}
\theoremstyle{remark}
\newtheorem*{remark}{Remark}

\begin{document}
\maketitle

\section*{Purpose}

The companion note \texttt{rmsea\_asymptotics.tex} fixes the meaning of RMSEA by
a Pitman drift: the misspecification is sent to zero at the $\sqrt n$ rate so the
test statistic converges to a (noncentral) $\chi^2$ and the noise inflation is
the first-order Satorra--Bentler trace $\tr(U\Gamma)$. This note drops the drift.
It expands the fitted discrepancy $\hat F$ to order $n^{-1}$ under
\emph{fixed} misspecification in both the covariance structure ($\sigma_0$ off
the model) and the distribution ($\Gamma\ne$ normal theory), for a general
minimum-discrepancy estimator. The point is to see what the local sequence
silently discards and to ask whether GLS-RMSEA differs from normal-theory RMSEA
at all once the drift is removed.

The single result is the bias expansion
\begin{equation}\label{eq:headline}
  \E[\hat F] \;=\; F_0 \;+\; \frac1n\Big[\;\underbrace{\tr(\tilde U\Gamma)}_{\text{generalized }df}
  \;+\; \underbrace{\beta_W}_{\text{weight influence, }\propto e}\;\Big] \;+\; o(n^{-1}),
\end{equation}
with two objects that the drift kills. The trace carries the
\emph{curvature-corrected} residual-weight matrix
$\tilde U = V - V\dot\sigma B^{-1}\dot\sigma'V$, built from the population
\emph{Hessian} $B=A-K$ of the discrepancy rather than its Gauss--Newton part
$A=\dot\sigma'V\dot\sigma$; the two agree only when the residual--curvature
coupling $K$ vanishes (correct specification, or a model linear in $\theta$). The
remainder $\beta_W$ is the influence of an estimated weight; it is identically
zero for a fixed weight (ULS) and vanishes as the residual $e\to0$ for any
estimator. Three corollaries read off \eqref{eq:headline}: a
curvature-corrected, robust bias correction for RMSEA
(Corollary~\ref{cor:bias}); an honest Wald confidence interval that needs no
noncentral-$\chi^2$ inversion (Corollary~\ref{cor:ci}); and the ML/GLS/ULS
comparison (Section~\ref{sec:estimators}).

Two structural remarks orient the rest. First, $\tilde U$ uses the Hessian $B$
where the first-order projector $U$ uses $A$; this is the fit-index sibling of
the observed-Hessian bread that replaces Fisher information in the
misspecification-robust standard errors of \texttt{misspec\_observed\_bread.tex},
and $\beta_W$ is the fit-index sibling of the data-dependent-weight meat term
there. Both corrections are zero at $H_0$ and leading-order off it; the
order-promotion is the same theorem twice. Second, in the language of
\texttt{aic\_tic\_rmsea\_bias.tex}, $\tr(\tilde U\Gamma)$ is to $\tr(U\Gamma)$ as
TIC's $\tr(I^{-1}J)$ is to AIC's $k$, and then one order deeper: the curvature
term is a correction past even the first-order robust (TIC-like) trace.

\section{Setup}\label{sec:setup}

Stack the $p^\ast=p(p{+}1)/2$ distinct sample covariances as $s=\vech\hat\Sigma$
(append means for a mean structure; the algebra is unchanged), with population
limit $s\pto\sigma_0$ and
\begin{equation}\label{eq:clt}
  z := \sqrt n\,(s-\sigma_0)\dto N(0,\Gamma),
\end{equation}
$\Gamma$ the asymptotic covariance of the moments, normal-theory or not. The
model is a smooth map $\theta\mapsto\sigma(\theta)$, $\theta\in\mathbb R^t$,
$df=p^\ast-t$. A minimum-discrepancy estimator minimizes a quadratic form in the
residual $r(\theta)=s-\sigma(\theta)$,
\begin{equation}\label{eq:disc}
  \hat F = \min_\theta\; r(\theta)'\hat V\,r(\theta),
  \qquad \hat V\pto V,
\end{equation}
with weight $\hat V$ either fixed (ULS) or a smooth function of the data (GLS).
Maximum likelihood is not exactly of this form; it is treated as the reweighted
limit in Section~\ref{sec:estimators}.

\paragraph{Pseudo-true value and population residual.}
Let
\begin{equation}\label{eq:pseudotrue}
  \theta_\ast = \argmin_\theta\,(\sigma_0-\sigma(\theta))'V(\sigma_0-\sigma(\theta)),
  \qquad e := \sigma_0-\sigma(\theta_\ast),\qquad F_0=e'Ve.
\end{equation}
Write $\dot\sigma=\partial\sigma/\partial\theta'$ at $\theta_\ast$
($p^\ast\times t$, full column rank) and the model Hessian tensor
$\ddot\sigma$, $(\ddot\sigma)_{a,ij}=\partial^2\sigma_a/\partial\theta_i\partial\theta_j$.
The first-order condition for \eqref{eq:pseudotrue} is the $V$-orthogonality of
the residual to the tangent space,
\begin{equation}\label{eq:popfoc}
  \dot\sigma'V e = 0 .
\end{equation}
Under correct specification $e=0$; under fixed misspecification $e=O(1)$ and
\eqref{eq:popfoc} is all that survives of it.

\paragraph{Curvature objects.}
Define the Gauss--Newton matrix, the residual--curvature coupling, and the
population Hessian of the discrepancy,
\begin{equation}\label{eq:ABK}
  A := \dot\sigma'V\dot\sigma,
  \qquad
  K_{ij} := (Ve)'(\ddot\sigma)_{:,ij}=\sum_a (Ve)_a(\ddot\sigma)_{a,ij},
  \qquad
  B := A-K = \tfrac12\,\partial^2_{\theta\theta}\big(e'Ve\big).
\end{equation}
$K$ is the only object in \eqref{eq:ABK} that requires both a nonzero residual
($Ve$) and a curved model ($\ddot\sigma\ne0$); it is the misspecification
signature. With $B$ in hand define the standard residual projector and its
curvature-corrected analogue
\begin{equation}\label{eq:Utilde}
  U := V-V\dot\sigma A^{-1}\dot\sigma'V,
  \qquad
  \tilde U := V-V\dot\sigma B^{-1}\dot\sigma'V = V-VPV,\quad P:=\dot\sigma B^{-1}\dot\sigma'.
\end{equation}
$U$ is a genuine residual projector ($U\dot\sigma=0$, idempotent in the
$V$-metric); $\tilde U$ is a perturbation of it that no longer annihilates the
tangent ($\tilde U\dot\sigma=-V\dot\sigma B^{-1}K$), reducing to $U$ as $K\to0$.
$U$ is the object of every limit law in the drift regime
(\texttt{rmsea\_asymptotics.tex}); $\tilde U$ is its fixed-misspecification
replacement.

\paragraph{Weight fluctuation.}
For an estimated weight expand
\begin{equation}\label{eq:weightexp}
  \hat V = V + n^{-1/2}\dot V + \tfrac12 n^{-1}\ddot V + o_p(n^{-1}),
\end{equation}
where $\dot V=O_p(1)$ is linear in $z$ (the derivative of the weight map along
the moment fluctuation) and $\ddot V=O_p(1)$ is the second-order term. Write the
residual-projected weight fluctuation
\begin{equation}\label{eq:w}
  w := \dot V\,e,
\end{equation}
an $O_p(1)$ vector that is identically zero for a fixed weight and is $O(\|e\|)$
otherwise. For GLS, $\dot V$ is driven by the raw moment fluctuation $z$; for ML,
by the projected fluctuation $\dot\sigma\,(\hat\theta-\theta_\ast)$
(Section~\ref{sec:estimators}).

\section{The stochastic expansion}

\begin{theorem}[Second-order expansion of $\hat F$]\label{thm:main}
Assume \eqref{eq:clt}, identification ($\dot\sigma$ full rank, $B$ nonsingular),
a weight that does not depend on $\theta$ (ULS, GLS, or any fixed/data weight),
and the usual smoothness. Then
\begin{equation}\label{eq:expansion}
  \hat F = F_0
  + n^{-1/2}\,\big(2e'Vz + e'\dot V e\big)
  + n^{-1}\,\Phi
  + o_p(n^{-1}),
\end{equation}
with
\begin{equation}\label{eq:Phi}
  \Phi = z'\tilde U z
       \;+\; 2\,w'(I-PV)z
       \;-\; w'Pw
       \;+\; \tfrac12\,e'\ddot V e .
\end{equation}
\end{theorem}

\begin{proof}
Write $u=\hat\theta-\theta_\ast=O_p(n^{-1/2})$ and split it as
$u=n^{-1/2}\xi+O_p(n^{-1})$. Expand the residual at the estimator,
\[
  r(\hat\theta)=e+\rho_1+\rho_2+o_p(n^{-1}),\quad
  \rho_1=n^{-1/2}z-\dot\sigma\,u,\quad
  \rho_2=-\tfrac12\,(\ddot\sigma)[u,u],
\]
where $(\ddot\sigma)[u,u]$ has $a$-component $u'(\ddot\sigma)_{:,a}u$. The
estimating equation $\dot\sigma(\hat\theta)'\hat V r(\hat\theta)=0$, expanded to
$O_p(n^{-1/2})$ with $\dot\sigma(\hat\theta)=\dot\sigma+(\ddot\sigma)[u,\cdot]$
and \eqref{eq:weightexp}, gives after using \eqref{eq:popfoc}
\[
  -A\,u + (\ddot\sigma)[u,\cdot]'Ve + n^{-1/2}\dot\sigma'Vz + n^{-1/2}\dot\sigma'\dot V e = o_p(n^{-1/2}).
\]
The second term is $Ku$ by \eqref{eq:ABK}; collecting,
\begin{equation}\label{eq:xi}
  (A-K)\,u = n^{-1/2}\dot\sigma'(Vz+\dot V e)+o_p(n^{-1/2})
  \;\Longrightarrow\;
  \xi = B^{-1}\dot\sigma'(Vz+w),\quad u=n^{-1/2}\xi.
\end{equation}
The weight fluctuation enters $\xi$ only through $w=\dot V e$, hence only under
misspecification. Now substitute into
$\hat F=(e+\rho)'\hat V(e+\rho)$, $\rho=\rho_1+\rho_2$, and use
\eqref{eq:popfoc} twice. The cross term with $\rho_1$ collapses to the linear
part,
$2e'V\rho_1=2n^{-1/2}e'Vz$ (the $\dot\sigma u_1$ piece dies by
\eqref{eq:popfoc}); the cross term with $\rho_2$ kills the second-order
parameter increment and leaves only curvature,
$2e'V\rho_2=-u'Ku$; and $e'\hat V e$ contributes the weight terms
$n^{-1/2}e'\dot V e+\tfrac12 n^{-1}e'\ddot V e$. Together with
$\rho_1'V\rho_1=n^{-1}(z'Vz-2\xi'\dot\sigma'Vz+\xi'A\xi)$ and the
$\dot V$--$\rho_1$ cross term $2n^{-1}\,w'\!\big(z-\dot\sigma\xi\big)$, the
$O(n^{-1})$ part is
\[
  \Phi = z'Vz - 2\xi'\dot\sigma'Vz + \xi'B\xi + 2w'(z-\dot\sigma\xi) + \tfrac12 e'\ddot V e,
\]
where $\xi'A\xi-u'Ku/n^{-1}$ combined to $\xi'B\xi$. Insert
$\xi=B^{-1}\dot\sigma'(Vz+w)$. With $P=\dot\sigma B^{-1}\dot\sigma'$ one gets
$\xi'B\xi=(Vz+w)'P(Vz+w)$, $\xi'\dot\sigma'Vz=(Vz+w)'PVz$ and
$\dot\sigma\xi=P(Vz+w)$. Expanding and cancelling,
\[
  \Phi = z'(V-VPV)z + 2w'(I-PV)z - w'Pw + \tfrac12 e'\ddot V e,
\]
which is \eqref{eq:Phi} with $\tilde U=V-VPV$.
\end{proof}

\begin{remark}[The cancellation is the content]
Everything hinges on the population FOC \eqref{eq:popfoc}. It removes the
tangent-direction sampling noise from the linear term (leaving $2e'Vz=2e'Uz$,
since $e'V=e'U$ by \eqref{eq:popfoc}) and removes the second-order parameter
increment $u_2$ from the bias, so the bias needs $\xi$ only to first order. What
remains is genuinely a residual-space object, projected by $\tilde U$.
\end{remark}

\section{Bias and the curvature-corrected df}

\begin{corollary}[Generalized degrees of freedom]\label{cor:bias}
Under the conditions of Theorem~\ref{thm:main},
\begin{equation}\label{eq:bias}
  \E[\hat F]=F_0+\frac1n\big[\tr(\tilde U\Gamma)+\beta_W\big]+o(n^{-1}),
  \qquad
  \beta_W=2\E\!\big[w'(I-PV)z\big]-\E\!\big[w'Pw\big]+\tfrac12 e'\E[\ddot V]\,e .
\end{equation}
For a fixed weight $\beta_W\equiv0$; for any estimator $\beta_W=O(\|e\|)$, so it
vanishes at $H_0$. Under correct specification $e=0$, $K=0$, $\tilde U=U$ and
\eqref{eq:bias} reduces to $\E[\hat F]=\tr(U\Gamma)/n+o(n^{-1})$, the
Satorra--Bentler mean $c\,df=\tr(U\Gamma)$ of \texttt{rmsea\_asymptotics.tex},
which is $df$ when $V=\Gamma^{-1}$.
\end{corollary}

\begin{proof}
Take expectations in \eqref{eq:expansion}. The $n^{-1/2}$ term has mean zero
($\E z=0$, $\E\dot V=0$). For the $n^{-1}$ term, $\E[z'\tilde Uz]=\tr(\tilde
U\Gamma)$ by the quadratic-form mean and the remaining terms define $\beta_W$.
\end{proof}

The estimand correction therefore generalizes the chain
\[
  df \;\longrightarrow\; \tr(U\Gamma)\;\longrightarrow\;\tr(\tilde U\Gamma)+\beta_W,
\]
in increasing order of fidelity: $df$ assumes a correctly weighted, correctly
distributed null; $\tr(U\Gamma)$ allows distributional misspecification but
freezes the model at its tangent; $\tr(\tilde U\Gamma)+\beta_W$ allows fixed
structural misspecification, correcting the projector for curvature and adding
the weight influence. A curvature-and-weight corrected RMSEA is
\begin{equation}\label{eq:rmseacorr}
  \rmsea_{\mathrm{ho}}
  = \sqrt{\max\!\Big(\frac{\hat F}{df}
      - \frac{\tr(\hat{\tilde U}\hat\Gamma)+\hat\beta_W}{n\,df},\,0\Big)} .
\end{equation}

\begin{remark}[Curvature is real for SEM]
$K\ne0$ requires $\ddot\sigma\ne0$. Factor and path models are nonlinear in
$\theta$ (loadings multiply, paths compose), so $\ddot\sigma\ne0$ generically and
the curvature term is present whenever the model is wrong. It vanishes for models
linear in $\theta$ (e.g.\ a free covariance matrix with linear equality
constraints), where $\tilde U=U$ exactly even under misspecification.
\end{remark}

\section{Comparison with the first-order trace}

The fixed-misspecification trace is the Satorra--Bentler trace plus an additive
series, so it slots directly against earlier robust RMSEA.

\begin{corollary}[Bread decomposition]\label{cor:decomp}
Let $N=\dot\sigma'V\Gamma V\dot\sigma$ (the parameter-space meat). Both traces are
the same total noise $\tr(V\Gamma)$ minus a model-absorbed part that differs only
in its bread:
\begin{equation}\label{eq:breadforms}
  \tr(U\Gamma)=\tr(V\Gamma)-\tr(A^{-1}N),
  \qquad
  \tr(\tilde U\Gamma)=\tr(V\Gamma)-\tr(B^{-1}N).
\end{equation}
The first-order (Satorra--Bentler) trace uses the Gauss--Newton bread $A^{-1}$;
the fixed-misspecification trace uses the Hessian bread $B^{-1}$. This is the same
$A^{-1}\!\to\!B^{-1}$ swap (expected information to observed Hessian) that turns a
naive standard error into a misspecification-robust one
(\texttt{misspec\_observed\_bread.tex}), here inside the fit-index trace.
Expanding $B^{-1}$ about $A^{-1}$ with $B=A-K$ and
$B^{-1}-A^{-1}=\sum_{m\ge1}(A^{-1}K)^mA^{-1}$ gives the additive series
\begin{equation}\label{eq:decompseries}
  \tr(\tilde U\Gamma)=\underbrace{\tr(U\Gamma)}_{\text{robust RMSEA}}
  \;-\;\tr(KM)\;-\;\tr(KA^{-1}KM)\;-\cdots,
  \qquad M:=A^{-1}NA^{-1},
\end{equation}
geometric in $A^{-1}K$ (ratio its spectral radius). The leading correction
$\mathrm{corr}_1=-\tr(KM)$ is $O(\|e\|)$: the residual-weighted model curvature
$K=(Ve)'\ddot\sigma$ contracted against $M$, the Gauss--Newton sandwich
covariance of $\hat\theta$ (the parameter covariance read at the tangent; it
equals the fixed-misspecification $\mathrm{acov}\,\hat\theta=B^{-1}NB^{-1}$ only
at $K=0$). Earlier robust RMSEA is the $K=0$ truncation.
\end{corollary}

\begin{remark}[Magnitudes, and ML]
In the ULS check (Section~\ref{sec:check}) the baseline $\tr(U\Gamma)=3.38$ and
the corrections $\mathrm{corr}_1=-0.55$, $\mathrm{corr}_2=-0.38$,
$\mathrm{corr}_3=-0.12,\dots$ sum to $\tr(\tilde U\Gamma)=2.26$; $\mathrm{corr}_1$
alone recovers half the gap, $\mathrm{corr}_{1:2}$ four-fifths. For ML there are
two baselines to deform from, the metric ($W_\ast\!\to\!V_0$) and the bread
($B\!\to\!A$), so the analogous expansion carries two families of corrections,
both $O(\|e\|)$ and both vanishing under correct specification.
\end{remark}

\section{The honest interval}

\begin{corollary}[CLT and Wald interval, no noncentral inversion]\label{cor:ci}
If $F_0>0$, then from \eqref{eq:expansion}
\begin{equation}\label{eq:clt2}
  \sqrt n\,(\hat F-F_0)\dto N(0,\omega^2),
  \qquad
  \omega^2=\Var\!\big(2e'Vz+e'\dot V e\big)=(2Ve+v_e)'\Gamma(2Ve+v_e),
\end{equation}
where $v_e$ is the vector representing the linear functional $z\mapsto e'\dot V
e=v_e'z$; for a fixed weight $v_e=0$ and $\omega^2=4e'V\Gamma Ve=4e'U\Gamma Ue$.
By the delta method, with $\varepsilon=\sqrt{F_0/df}$,
\begin{equation}\label{eq:rmseaclt}
  \sqrt n\,(\rmsea-\varepsilon)\dto N\!\Big(0,\;\frac{\omega^2}{4\,df\,F_0}\Big),
\end{equation}
giving a Wald interval for $\varepsilon$ directly on the discrepancy scale.
\end{corollary}

Equation \eqref{eq:rmseaclt} is the fixed-misspecification interval: under a
genuinely wrong model the statistic is asymptotically normal, not noncentral
$\chi^2$, and the interval is a plain sandwich, not a CDF inversion. It
complements rather than replaces the noncentral interval of
\texttt{rmsea\_asymptotics.tex}: that one is built for the near-null regime
($F_0\downarrow0$) where \eqref{eq:rmseaclt} degenerates ($\omega^2/F_0$ blows
up), while \eqref{eq:rmseaclt} is built for clearly-misspecified models where the
noncentral inversion is least trustworthy. The crossover is exactly where
neither device is comfortable, which is the honest statement of where RMSEA
interval coverage is hard.

\section{What the drift discards}

Put back the Pitman sequence $e_n=n^{-1/2}\delta$ with $\delta$ a fixed off-model
direction, the regime of \texttt{rmsea\_asymptotics.tex}. Then $Ve_n=O(n^{-1/2})$
so $K=O(n^{-1/2})\to0$, hence $B\to A$ and $\tilde U\to U$; the weight influence
$w=\dot V e_n=O_p(n^{-1/2})\to0$, so $\beta_W\to0$; and the normal CLT
\eqref{eq:clt2} degenerates while $n\hat F$ acquires the (noncentral) $\chi^2$
limit. The drift annihilates precisely the two corrections this note is about.

\begin{center}
\begin{tabular}{p{0.27\linewidth}p{0.32\linewidth}p{0.31\linewidth}}
 & fixed misspecification ($e=O(1)$) & Pitman drift ($e=n^{-1/2}\delta$)\\
\hline
limit law of $\hat F$ & normal, $\sqrt n(\hat F-F_0)\dto N(0,\omega^2)$ &
$n\hat F\dto$ (noncentral) $\chi^2$ mixture\\
noise/bias mean & $\tr(\tilde U\Gamma)$, curvature-corrected &
$\tr(U\Gamma)=c\,df$, no curvature\\
weight influence $\beta_W$ & $O(\|e\|)$, present & $\to0$\\
ML vs GLS & differ (two-metric, pseudo-true) & coincide to leading order\\
interval & Wald via $\omega^2$ \eqref{eq:rmseaclt} & noncentral-$\chi^2$ inversion\\
\end{tabular}
\end{center}

The local-asymptotics motivation is not wrong; it is a deliberate choice to work
in the only regime where the $df$ subtraction is exact and the curvature and
weight terms are negligible. Outside it, those terms are the leading description
of a wrong model's fit.

\section{ML, GLS, ULS}\label{sec:estimators}

The estimators differ in where the weight is anchored, and once the drift is gone
this separates their RMSEAs. The unifying tool is that $\tilde U$ is, for
\emph{any} smooth discrepancy, half the Hessian of the profiled criterion.

\begin{lemma}[Profiled-Hessian form]\label{lem:profhess}
For a smooth discrepancy $F(s,\theta)$ with population minimizer $\theta_\ast$,
Corollary~\ref{cor:bias} holds with
\begin{equation}\label{eq:profhess}
  \tilde U \;=\; \tfrac12\,\nabla_s^2\Big[\min_\theta F(s,\theta)\Big]\Big|_{s=\sigma_0}
  \;=\; V_{ss}-\tfrac14\,C\,B^{-1}C',
\end{equation}
where $V_{ss}=\tfrac12\partial^2_{ss}F$, $C=\partial^2_{s\theta}F$ and
$B=\tfrac12\partial^2_{\theta\theta}F$ at $(\sigma_0,\theta_\ast)$. The
fixed-weight quadratic \eqref{eq:Utilde} is the case $V_{ss}=V$,
$C=-2V\dot\sigma$ (then $\tfrac14 CB^{-1}C'=V\dot\sigma B^{-1}\dot\sigma'V$). A
non-quadratic $F$ enters only through these three second derivatives; there is no
separate higher-derivative term at $O(1/n)$, and the data skewness enters the
bias only at $O(n^{-3/2})$. For a data-dependent weight $V(s)$ the $s$-derivatives
of the weight populate $V_{ss}$ and $C$, reproducing $\beta_W$ of
Corollary~\ref{cor:bias} in trace form.
\end{lemma}

\paragraph{ULS.} Fixed weight, $\dot V=0$, so $\beta_W\equiv0$ and the bias is
purely $\tr(\tilde U\Gamma)/n$. ULS is never efficient, so $\tr(U\Gamma)\ne df$
already at leading order regardless of normality; the curvature term is then a
clean second correction with no weight bookkeeping. This is the cleanest place to
exhibit $\tilde U\ne U$.

\paragraph{GLS.} Weight $\hat V=V(s)$ at the sample moments, so $\dot V$ is the
derivative of $V(\cdot)$ along the \emph{raw} fluctuation $z$. The full $\beta_W$
is present, leading-order in the residual through the cross term
$2\E[w'(I-PV)z]$, and $\omega^2$ in \eqref{eq:clt2} acquires the $v_e$ term.

\paragraph{ML.} ML is not a fixed quadratic: its normal equation
$\dot\sigma'W(\theta)\,(\sigma_0-\sigma(\theta))=0$ anchors the weight
$W(\theta)=\tfrac12 D'(\Sigma(\theta)^{-1}\otimes\Sigma(\theta)^{-1})D$ at the
\emph{fitted model}. Two consequences. First, the pseudo-true value solves a
model-metric projection, so under fixed misspecification
$\theta_\ast^{\mathrm{ML}}\ne\theta_\ast^{\mathrm{GLS}}$ and
$F_0^{\mathrm{ML}}\ne F_0^{\mathrm{GLS}}$: the two estimate \emph{different
population numbers}, already at zeroth order, measuring misfit in different
metrics (Olsson, Foss, Troye and Howell \cite{olssonetal2000}). Second, by
Lemma~\ref{lem:profhess} the noise weight is the moment-Hessian
$V_{ss}=V_0=\tfrac12 D'(\Sigma_0^{-1}\otimes\Sigma_0^{-1})D$ at the \emph{data},
while the cross-term $C=-2W_\ast\dot\sigma$ carries the \emph{model} weight
$W_\ast=\tfrac12 D'(\Sigma_\ast^{-1}\otimes\Sigma_\ast^{-1})D$,
$\Sigma_\ast=\Sigma(\theta_\ast)$, so
\begin{equation}\label{eq:UtildeML}
  \tilde U_{\mathrm{ML}} = V_0 - W_\ast\dot\sigma\,B_{\mathrm{ML}}^{-1}\dot\sigma'W_\ast .
\end{equation}
This is the two-metric object: sampling noise weighted in the data metric,
residuals projected in the model metric. The non-quadraticity of the log-det is
absorbed entirely into $(V_0,W_\ast,B_{\mathrm{ML}})$; there is no separate
third-derivative term (correcting an earlier draft). It all collapses under
correct specification ($\Sigma_\ast=\Sigma_0$, $W_\ast=V_0$, $e=0$), where ML, GLS
and ULS coincide to the order that matters. The numerical check confirms
\eqref{eq:UtildeML} to finite-difference precision.

So the answer to ``does GLS-RMSEA do anything that normal-theory RMSEA does
not'' is yes, and sharper than expected: under fixed misspecification the two
already target different estimands (different metric), and on top of that carry
different $O(1/n)$ biases through $\beta_W$ and $B$. Under the drift, all of this
disappears and the classical equivalence is restored. The daylight between
estimators lives entirely in the fixed-misspecification, higher-order regime.

\section{Numerical check}\label{sec:check}

A one-factor model is fit to a two-factor population (six indicators, three per
factor, factor correlation $0.55$), a misspecification the one-factor model
cannot absorb ($df=9$; ULS $F_0=0.23$, ML $F_0=0.49$, the differing estimands of
Section~\ref{sec:estimators}). Two checks.

\emph{Deterministic.} By Lemma~\ref{lem:profhess} the bias coefficient is half
the Hessian of the profiled criterion, which finite differencing pins exactly,
with no Monte Carlo. Comparing $\tfrac12\nabla_s^2[\min_\theta F]$ at $\sigma_0$
to the closed forms:
\begin{center}
\begin{tabular}{lccc}
 & $\|\tfrac12 H_{\mathrm{prof}}-\tilde U\|/\|\tilde U\|$ & $\tr(\tilde U\Gamma_{\mathrm{NT}})$ corrected & naive (Gauss--Newton $A$)\\
\hline
ULS & $2\times10^{-8}$ & $2.26$ & $3.38$\\
ML  & $9\times10^{-6}$ & $1.24$ & $9.07$\\
\end{tabular}
\end{center}
ULS is a single-metric quadratic (positive control); ML is the two-metric
$\tilde U_{\mathrm{ML}}$ of \eqref{eq:UtildeML}, confirmed to finite-difference
precision. Replacing the Hessian $B$ by Gauss--Newton $A$ is wrong by a factor of
several, for ML by $7\times$ (and with the opposite conclusion about how much
sampling noise a wrong model spends): discarding the curvature is not a small
error.

\emph{Monte Carlo (ULS, $\beta_W=0$).} Confirms the Hessian-to-bias step,
$N(\E[\hat F]-F_0)\to\tr(\tilde U\Gamma)$, under normal and $t_6$ data; the sample
covariance is the unbiased $\div(N{-}1)$ form (the $\div N$ form carries its own
$O(1/N)$ moment bias contaminating this order), $\Gamma_{\mathrm{NT}}$ closed form
under normality and Monte-Carlo under $t_6$:
\begin{center}
\begin{tabular}{llccc}
data & $N$ & $N(\E[\hat F]-F_0)$ (MC se) & naive $\tr(U\Gamma)$ & corrected $\tr(\tilde U\Gamma)$\\
\hline
normal & $200$ & $2.23\ (0.05)$ & $3.38$ & $\mathbf{2.26}$\\
normal & $800$ & $2.21\ (0.10)$ & $3.38$ & $\mathbf{2.26}$\\
$t_6$  & $200$ & $4.22\ (0.08)$ & $7.00$ & $\mathbf{4.74}$\\
$t_6$  & $800$ & $4.78\ (0.16)$ & $7.00$ & $\mathbf{4.74}$\\
\end{tabular}
\end{center}
By $N=800$ the empirical bias is within one Monte-Carlo standard error of
$\tr(\tilde U\Gamma)$ in both conditions; the naive trace is off by the full gap
(over twenty standard errors at $N=200$, normal). ML is validated
deterministically rather than by Monte Carlo: its bias has a large higher-order
tail and small-$N$ improper-solution inflation that the exact Hessian check
sidesteps. The curvature term is not a third-decimal refinement; it is a third to
a half of the whole bias. Script:
\texttt{higher\_order\_discrepancy\_check.py}.

\section{Connections and caveats}

\begin{remark}[Hessian bread, weight meat: the SE sibling]
$\tilde U$ uses the population Hessian $B=A-K$ where $U$ uses the Gauss--Newton
$A$; under correct specification $K=0$ and they agree, the discrepancy-function
analogue of the information equality. This is the same substitution that the
misspecification-robust standard errors of \texttt{misspec\_observed\_bread.tex}
make when they replace the Fisher bread by the observed Hessian, and $\beta_W$ is
the fit-index counterpart of their data-dependent-weight meat term. Both
corrections are $O(\|e\|)$, hence zero at $H_0$ and leading-order off it: the
order-promotion (Newey--McFadden null irrelevance, Hall--Inoue off-null
promotion) is one phenomenon appearing in the variance there and in the bias
here.
\end{remark}

\begin{remark}[TIC analogy, one order deeper]
\texttt{aic\_tic\_rmsea\_bias.tex} reads $df$ as the AIC penalty and notes the
robust trace as its TIC analogue. Here $\tr(U\Gamma)$ is that TIC-like trace
(first order), and $\tr(\tilde U\Gamma)+\beta_W$ is the next order. The user's
instinct stands: the $df$ correction is AIC-like (a fixed integer), the robust
trace is TIC-like (a sandwich trace), and the practical complaint about TIC, that
$\tr(\hat I^{-1}\hat J)$ is noisy because $\hat J$ is a fourth-moment object,
transfers verbatim: $\tr(\hat{\tilde U}\hat\Gamma)$ inherits the sampling
variance of $\hat\Gamma$. The higher-order expansion is what lets one judge the
trade: when $\|e\|$ and the kurtosis are small the AIC-like $df$ is nearly
unbiased and lower variance; when they are not, the bias of the cheap correction
is exactly \eqref{eq:bias} and may dominate. A structured (model-implied or
shrunk) $\hat\Gamma$ is the natural middle estimator, the RMSEA analogue of a
regularized TIC.
\end{remark}

\begin{remark}[Caveats]
The truncation $\max(\cdot,0)$ in \eqref{eq:rmseacorr} re-introduces an
$O(n^{-1/2})$ bias near $F_0=0$, where the fixed-misspecification expansion is
least relevant and the near-null device should take over. Data skewness enters
the bias only at $O(n^{-3/2})$ (Lemma~\ref{lem:profhess}); the $t_6$ shortfall at
$N=200$ in the table is that tail. The expansion assumes finite fourth moments
for $\Gamma$ and finite sixth moments for the $n^{-1}$ bias to have finite
expectation; under heavier tails \eqref{eq:bias} is a stochastic-order statement,
not a moment statement.
\end{remark}

\begin{thebibliography}{9}

\bibitem{browne1984}
Browne, M. W. (1984).
Asymptotically distribution-free methods for the analysis of covariance
structures.
\emph{British Journal of Mathematical and Statistical Psychology}, 37(1),
62--83.
\href{https://doi.org/10.1111/j.2044-8317.1984.tb00789.x}{doi:10.1111/j.2044-8317.1984.tb00789.x}.

\bibitem{olssonetal2000}
Olsson, U. H., Foss, T., Troye, S. V., \& Howell, R. D. (2000).
The performance of ML, GLS, and WLS estimation in structural equation modeling
under conditions of misspecification and nonnormality.
\emph{Structural Equation Modeling}, 7(4), 557--595.
\href{https://doi.org/10.1207/S15328007SEM0704_3}{doi:10.1207/S15328007SEM0704\_3}.

\bibitem{satorra1989}
Satorra, A. (1989).
Alternative test criteria in covariance structure analysis: A unified approach.
\emph{Psychometrika}, 54(1), 131--151.
\href{https://doi.org/10.1007/BF02294453}{doi:10.1007/BF02294453}.

\bibitem{satorrabentler1994}
Satorra, A., \& Bentler, P. M. (1994).
Corrections to test statistics and standard errors in covariance structure
analysis. In A. von Eye \& C. C. Clogg (Eds.), \emph{Latent Variables Analysis}
(pp.\ 399--419). Sage.

\bibitem{shapiro1983}
Shapiro, A. (1983).
Asymptotic distribution theory in the analysis of covariance structures (a
unified approach).
\emph{South African Statistical Journal}, 17(1), 33--81.

\bibitem{shapiro2007}
Shapiro, A. (2007).
Statistical inference of moment structures. In S.-Y. Lee (Ed.),
\emph{Handbook of Latent Variable and Related Models} (pp.\ 229--260). Elsevier.
\href{https://doi.org/10.1016/B978-044452044-9/50013-9}{doi:10.1016/B978-044452044-9/50013-9}.

\bibitem{white1982}
White, H. (1982).
Maximum likelihood estimation of misspecified models.
\emph{Econometrica}, 50(1), 1--25.
\href{https://doi.org/10.2307/1912526}{doi:10.2307/1912526}.

\bibitem{yuanhayashibentler2007}
Yuan, K.-H., Hayashi, K., \& Bentler, P. M. (2007).
Normal theory likelihood ratio statistic for mean and covariance structure
analysis under alternative hypotheses.
\emph{Journal of Multivariate Analysis}, 98(6), 1262--1282.
\href{https://doi.org/10.1016/j.jmva.2006.08.005}{doi:10.1016/j.jmva.2006.08.005}.

\end{thebibliography}

\end{document}
